\uuid{7ol1}
\exo7id{3652}
\titre{exo7 3652}
\auteur{quercia}
\organisation{exo7}
\datecreate{2010-03-10}
\isIndication{false}
\isCorrection{false}
\chapitre{Espace vectoriel}
\sousChapitre{Somme directe}
\cpge{mpsi-mp2i.espaces-vectoriels.d, pcsi.espaces-vectoriels.d, ptsi.espaces-vectoriels.d, mpsi-mp2i.dimension-finie.c, pcsi.dimension-finie.c, ptsi.dimension-finie.c, mpsi-mp2i.polynomes-fractions-rationnelles.c, pcsi.polynomes.c, ptsi.polynomes.c, bcpst1.polynomes-reels.b}
\module{Algèbre}
\niveau{L1}
\difficulte{}

\contenu{
\texte{
Soit $E =  K_3[X]$,
$F = \{ P \in E$ tq $P(0) = P(1) = P(2) = 0 \}$,
$G = \{ P \in E$ tq $P(1) = P(2) = P(3) = 0 \}$,
et $H = \{ P \in E$ tq $P(X) = P(-X) \}$.
}
\begin{enumerate}
    \item \question{Montrer que $F \oplus G = \{ P \in E$ tq $P(1) = P(2) = 0 \}$.}
    \item \question{Montrer que $F \oplus G \oplus H = E$.}
\end{enumerate}
}
